\documentclass[10pt,twocolumn,letterpaper]{article}

\usepackage{wacv}
% (tắt cảnh báo của caption)
\usepackage{silence}
\WarningFilter{caption}{Unknown document class}
\usepackage{times}
\usepackage{epsfig}
\usepackage{graphicx}
\usepackage{amsmath}
\usepackage{amssymb}
\usepackage{booktabs}
\usepackage{array}
\usepackage{algorithm}
\usepackage{algpseudocode}
\usepackage{svg}
\usepackage[none]{hyphenat}
\usepackage{enumitem}
\usepackage{multirow}  
\usepackage{changepage}
\usepackage[compatibility=false]{caption}
\usepackage{float}
\captionsetup[figure]{justification=justified}
\captionsetup[table]{justification=justified}
\raggedbottom  

\usepackage[pagebackref=true,breaklinks=true,colorlinks=false,bookmarks=false]{hyperref}

%\linespread{1.5}\selectfont % Set line spacing to 1.5
\wacvfinalcopy % *** Uncomment this line for the final submission

%\def\wacvPaperID{****} % *** Enter the wacv Paper ID here

\def\httilde{\mbox{\tt\raisebox{-.5ex}{\symbol{126}}}}

% Pages are numbered in submission mode, and unnumbered in camera-ready
\ifwacvfinal\pagestyle{empty}\fi
\setcounter{page}{1}

\setlength{\parindent}{0pt}


\makeatletter
\def\wacvsect#1{\wacvsection{#1}}
\def\wacvsubsect#1{\wacvsubsection{#1}}

% Custom title formatting for better prominence
\renewcommand{\@maketitle}{%
  \newpage
  \null
  \vskip 1.5em%
  \begin{center}%
  \let \footnote \thanks
    {\Large \bfseries \@title \par}%
    \vskip 1em%
    {\normalsize
      \lineskip .5em%
      \begin{tabular}[t]{c}%
        \@author
      \end{tabular}\par}%
    \vskip 0.5em%
    % {\normalsize \@date}%
  \end{center}%
  \par
  \vskip 1em}
\makeatother

\captionsetup{justification=centering}
\begin{document}

%%%%%%%%% TITLE
\title{Hybrid Search Architecture with Multi-Stage Fusion and Knowledge-Aware Prompt Engineering for Retrieval Augmented Generation}                     

% Authors at different institutions
\author{Luyen Nguyen Tien \\
{\tt\small bit220101@st.cmcu.edu.vn}
\\[1ex]
CMC University
\and
Binh Hoang Tieu $\footnote{Corresponding author}$\\
{\tt\small htbinh@cmcu.edu.vn}
\\[1ex]
CMC University
}


\maketitle
\ifwacvfinal\thispagestyle{empty}\fi 
%%%%%%%%% ABSTRACT
\begin{abstract}
\begin{adjustwidth}{0.02\textwidth}{0.02\textwidth}
  Retrieval-Augmented Generation (RAG) systems combine information retrieval with generative models to deliver accurate responses, yet many frameworks still face limitations in retrieval fusion and prompt optimization. We introduce HAPE-RAG with two key innovations: (1) a multi-stage hybrid search architecture that couples semantic vector search with POS/LLM-driven keywording, duplicate-aware fusion, and batched BGE re-ranking; and (2) a knowledge-aware prompting system that uses a structured, pattern-based template integrating system instruction, KB metadata, Top-$k$ context with document names, user query, instruction template, and conversation memory. Experiments on Vietnam's Undergraduate Training Regulations (VUTR) and Narrative Comprehension (NarrativeQA) datasets show consistent gains over semantic-only, BM25-only, and naive hybrids. On VUTR/NarrativeQA, we observe RAGAS Faithfulness of 0.928/0.890 and Answer Relevancy of 0.932/0.792, together with improved ranking NDCG to 0.875 and 0.862 respectively, with a 70--30 (semantic–keyword) fusion emerging as a robust setting. Compared to a strong flat-concatenation hybrid baseline, HAPE-RAG improves Faithfulness by \textbf{+2.09\%}, Answer Relevancy by \textbf{+2.98\%}, BLEU-1 by \textbf{+1.6} points, and Retrieval Accuracy by \textbf{+0.031} points (\textbf{+3.54\%}).
\end{adjustwidth}
\end{abstract}

%%%%%%%%% KEYWORDS
\begin{adjustwidth}{0.02\textwidth}{0.02\textwidth}
\noindent\textbf{Keywords:} \textit{RAG, Hybrid Search, Prompt Engineering, Natural Language Processing}
\end{adjustwidth}

%%%%%%%%% BODY TEXT
\section{Introduction}

Retrieval-Augmented Generation (RAG) has emerged as a transformative paradigm that bridges the gap between information retrieval and large language models (LLMs), enabling systems to generate accurate, contextually relevant responses by leveraging external knowledge sources ~\cite{lewis2020retrieval}. The evolution of RAG systems has progressed through several generations, from simple keyword-based retrieval to sophisticated neural approaches that leverage deep learning for semantic understanding  ~\cite{vaswani2017attention,devlin2019bert}.

\vspace{0.5em}
However, the increasing complexity and diversity of modern RAG frameworks have introduced two critical challenges that limit their practical effectiveness. First, existing retrieval approaches typically rely on single-modal methods, either semantic vector search or traditional keyword matching, which fail to capture the complementary strengths of both paradigms  ~\cite{gao2021retrieval,mao2023neural}. Semantic search excels at understanding conceptual relationships and contextual meaning, while keyword search provides precise matching for specific terms and entities. The lack of effective fusion strategies between these modalities results in suboptimal retrieval performance, with studies showing 15-25\% improvements in precision when hybrid approaches are properly implemented ~\cite{zhang2024semantic,liu2024impact}. A second bottleneck is prompting. Static, flat templates seldom exploit knowledge‑base (KB) metadata or conversation memory, causing field leakage and weaker attribution ~\cite{wei2022chain,liu2023chained}. In practice, RAG deployments must also manage latency and cost; unstructured prompts and one‑pass fusion exacerbate both. These two fundamental issues - the lack of effective hybrid retrieval strategies and the absence of knowledge-aware prompt engineering - represent the primary bottlenecks in modern RAG system development and deployment, limiting their applicability across diverse domains and use cases.

\vspace{0.5em}
 We introduce HAPE-RAG, a deployment‑oriented RAG framework that (i) uses a \emph{multi‑stage dense–sparse hybrid}-semantic search plus POS/LLM keyphrases, duplicate‑aware fusion, and neural reranking to improve top‑k evidence; (ii) adopts a \emph{knowledge‑aware structured prompt} (JSON‑style) that explicitly integrates system instruction, KB metadata, top‑k context with document names, query, instruction template, and conversation memory, strengthening grounding and attribution; and (iii) provides \emph{production engineering} (batching, singleton model management, caching, adaptive cutoffs) to balance quality and latency. Experiments on VUTR and NarrativeQA show consistent gains over semantic‑only, keyword‑only.

\section{Related Works}

\textbf{Hybrid Search.}
Classical work studied multiple‑evidence combination via linear weighting and voting schemes  ~\cite{lee1997combining}. Modern dense–sparse hybrids leverage complementary strengths: dense encoders offer paraphrase‑robust semantics (e.g., DPR/Contriever/ColBERT  ~\cite{karpukhin2020dpr,izacard2021contriever,khattab2020colbert}), while sparse/keyword methods (BM25/SPLADE  ~\cite{robertson2009bm25,formal2021splade}) recover rare terms, identifiers, and entities  ~\cite{gao2021retrieval}. Neural rerankers (cross‑encoders, BGE  ~\cite{nogueira2019rerankbert,xu2023bge}) substantially improve ordering but add latency/compute  ~\cite{mao2023neural}. However, many hybrids remain \emph{single‑pass}: they fuse scores without staged candidate expansion, duplicate control, or cost‑aware batching/caching. This limits recall–precision trade‑offs and propagates noisy, redundant passages. Our approach adopts a \emph{four‑stage} pipeline: semantic retrieval, POS/LLM keywording, duplicate‑aware fusion, and neural reranking—implemented with batching, singleton loading, and caches, aligning retrieval quality with production constraints.

\vspace{0.5em}
\textbf{Prompt Engineering.}
Template prompts and instruction tuning are widely used  ~\cite{wei2022chain}; dynamic/chained prompting improves reasoning in multi‑step tasks ~\cite{liu2023chained,wang2022selfconsistency,yao2023react}. Yet KB‑aware \emph{structuring}—explicit fields for system instruction, KB metadata, top‑$k$ context with document names, user query, instruction template, and conversation memory—remains underexplored for grounding and attribution at inference time. Flat concatenation blurs fields, increases leakage, and inflates tokens. We contribute a JSON‑style schema that keeps roles separable, supports principled token budgeting, and empirically raises faithfulness and answer relevancy.

\vspace{0.5em}
\textbf{RAG Architecture and Systems.}
Multi‑stage RAG improves over single‑method pipelines ~\cite{gao2021retrieval,lewis2020retrieval} and recent surveys/system papers highlight the importance of concurrency/caching and vector stores (e.g., FAISS/pgvector) for scale ~\cite{gao2021retrieval,johnson2019faiss,pgvector2021,zhang2024semantic,liu2024impact}. We extend this line by tightly coupling staged dense–sparse retrieval with neural reranking and a structured prompt layer, while exposing system levers (candidate cutoffs, batch sizes, cache TTLs) that make quality–cost trade‑offs explicit for deployment.

%------------------------------------------------------------------------
\begin{figure*}[!htbp]
    \centering
    \includesvg[width=\textwidth]{hape-rag.svg}
    \caption{HAPE-RAG pipeline. User query $q$ is clarified to $q'$ using conversation memory $\mathbf{M}$, then processed in parallel: (i) embedded to $\mathbf{e}$ for semantic pgvector search (Top-$p$, pink blocks), (ii) keyword extraction for array-overlap search (Top-$n$, blue blocks). Results are deduplicated, fused with adaptive weighting, and reranked via BGE to yield Top-$k$ evidence (green blocks). A structured JSON prompt $P$ combines system instructions, KB metadata $\mathbf{K}$, retrieved context, and $\mathbf{M}$; the LLM generates the answer and updates $\mathbf{M}$ for continuous learning.}
    \label{fig:system_architecture}
\end{figure*}

\section{Methodology}

We introduce HAPE-RAG, a novel approach for Retrieval Augmented Generation through two fundamental innovations: a multi-stage hybrid search architecture that synchronously combines semantic retrieval, retrieval, and re-ranking, and a knowledge-aware prompting system that integrates in a structured way to address the blind spot problem in keypoint recognition, with the architectural details shown in Figure \ref{fig:system_architecture}.

\subsection{Multi-Stage Hybrid Search Architecture}

In the first stage, we perform a dense embedding of the original query $q$ into a more detailed and complete sentence $q'$ based on the current context using LLM. Given the query $q'$, we generate a dense embedding vector $\mathbf{e}_q \in \mathbb{R}^d$ \quad via the embedding model, where $d$ represents the embedding dimension, and search for dense vector similarity using the pgvector extension in PostgreSQL. The semantic search engine will retrieve $k$ of the Top-$p$ text segments with the highest semantic similarity based on the cosine similarity calculation:
$$\text{sim}(q, c_i) = \frac{\mathbf{e}_q \cdot \mathbf{e}_{c_i}}{\lVert \mathbf{e}_q\rVert_2 \cdot \lVert\mathbf{e}_{c_i}\rVert_2}$$
where $\mathbf{e}_{c_i}$ denotes the embedding vector of segment $c_i$.

\vspace{0.5em}
The second phase deploys dual keyword extraction strategies to capture both the linguistic and semantic aspects of the query. The NLP-based approach uses a tokenizer library for Part of Speech (POS) tagging, focusing on noun categories. The extraction process applies linguistic filtering to remove stop words and shortened tokens. We also combine using LLM to generate 5-7 keyword clusters, each consisting of 2-4 words, focusing on specialized terminology and core conceptual entities. The extracted keywords are effectively used for array overlap search via the intersection operator:
$$\mathcal{K}{query} \cap \mathcal{K}{chunk} \neq \emptyset$$
where $\mathcal{K}_{query}$ represents the query keyword set and $\mathcal{K}_{chunk}$ denotes the chunk keyword array. Finally, the Top-$n$ most relevant chunks are retrieved.

\vspace{0.5em}
The fusion stage implements intelligent deduplication and score normalization to combine results from both retrieval modalities. The system maintains a comprehensive set of retrieved chunk identifiers to eliminate duplicates while preserving the diversity inherent in multi-stage retrieval. The fusion process employs an adaptive weighted combination strategy:
$$\text{score}(c_i) = \alpha \cdot \text{score}_{semantic}(c_i) + \beta \cdot \text{score}_{keyword}(c_i)$$

where $\alpha,\beta\in[0,1]$ and $\alpha+\beta=1$. We first min–max normalize scores from each modality so they are on a comparable scale, then fuse them with a validation‑tuned default of $\alpha=0.7$ (semantic) and $\beta=0.3$ (keyword). For queries dominated by exact terms/entities—indicated by high POS/NER density or strong query–chunk keyword overlap—we slightly increase $\beta$ to give the keyword branch more influence; for generic or paraphrastic queries we keep $\alpha$ higher to prioritize semantics. In practice this adaptive mix reduces misses on rare terms without sacrificing semantic robustness. The final stage uses \textit{BAAI/bge-reranker-base} for neural reranking.

\subsection{Prompt Engineering System}

The prompt engineering framework addresses the limitations of static template-based approaches through dynamic integration of knowledge base metadata and conversation context. The system employs a modular template construction methodology that enables fine-grained control over response generation while maintaining contextual awareness.

\vspace{0.5em}
The prompt construction follows a structured template system combining multiple information sources:
$$\mathcal{P} = \mathcal{S} \oplus \mathcal{M} \oplus \mathcal{K} \oplus \mathcal{C} \oplus \mathcal{Q} \oplus \mathcal{I}$$
where $\mathcal{S}$ represents \textit{S}ystem instruction, $\mathcal{M}$ denotes conversation \textit{M}emory, $\mathcal{K}$ signifies \textit{K}nowledge base metadata, $\mathcal{C}$ indicates \textit{C}ontext with document names, $\mathcal{Q}$ represents \textit{Q}uery section, and $\mathcal{I}$ denotes \textit{I}nstruction template.

\vspace{0.5em}
The system dynamically integrates knowledge base metadata through the \texttt{kb\textbackslash info} parameter, incorporating knowledge base titles, descriptions, and document metadata. The integration process employs template variable substitution with the following formulation:
$$\mathcal{K} = \text{Template}{KB} \cdot \text{Substitute}(\text{metadata}{KB})$$
where $\text{Template}_{KB}$ represents the knowledge base template structure and $\text{metadata}{KB}$ denotes the extracted metadata information.

\vspace{0.5em}
The conversation memory system implements a structured JSON schema that tracks multiple aspects of the dialogue:
$$\mathcal{M} = \{\text{summary}, \text{topics}, \text{key\_points}, \text{context}, \text{metadata}\}$$
The memory system stores the conversation summary, identified discussion topics, key information points, recent query history, current context focus, and metadata including message count and timestamp. This makes responses more coherent throughout the conversation.

\vspace{0.5em}
With the components in the structure, we aggregate them into JSON - a standard that matches LLM's Prompt, rather than concatenating them all into a single paragraph or semantic summary as is currently the case. In practice, a knowledge base may contain multiple documents, which can lead to confusion when handling large content volumes and queries that span fragmented sections across those documents. This confusion may prevent the LLM from locating the correct context, resulting in inaccurate or missing answers. To address this, we add the document name at the beginning of each fragment, along with a brief description of the knowledge to help the context be tight and work correctly when fed into LLM's answer. Figure \ref{fig:system_prompt} clearly shows our Prompt structure.

\begin{figure}[!ht]
    \centering
    \includesvg[width=.8\columnwidth,keepaspectratio]{prompt.svg}
    \caption{Structured prompt schema: (i) system instruction, (ii) KB metadata, (iii) Top-$k$ context with document names, (iv) query section, (v) instruction template, and (vi) conversation memory. This JSON structure improves grounding versus plain concatenation.}
    \label{fig:system_prompt}
\end{figure}

\begin{table*}[t]
\small
\centering
\caption{Performance across datasets. ``Ours'' = HAPE-RAG. Columns: \textbf{Sem-only} = Semantic-only RAG (dense retrieval), \textbf{BM25-only} = Keyword-only RAG (BM25), \textbf{Naive} = Naive Hybrid RAG (simple concat of Sem + BM25).}
\label{tab:overall_performance}
\setlength{\tabcolsep}{5.5pt}
\begin{tabular}{lcccc cccc}
\toprule
& \multicolumn{4}{c}{\textbf{VUTR}} & \multicolumn{4}{c}{\textbf{NarrativeQA}} \\
\cmidrule(lr){2-5}\cmidrule(lr){6-9}
\textbf{Metric} & \textbf{Ours} & Sem-only & BM25-only & Naive & \textbf{Ours} & Sem-only & BM25-only & Naive \\
\midrule
RAGAS Faithfulness      & \textbf{0.928} & 0.909 & 0.874 & 0.916 & \textbf{0.890} & 0.861 & 0.835 & 0.872 \\
RAGAS Answer Relevancy  & \textbf{0.932} & 0.905 & 0.866 & 0.912 & \textbf{0.792} & 0.763 & 0.744 & 0.776 \\
RAGAS Context Relevancy & \textbf{0.921} & 0.900 & 0.855 & 0.905 & \textbf{0.889} & 0.860 & 0.836 & 0.875 \\
RAGAS Context Recall    & \textbf{0.906} & 0.887 & 0.838 & 0.892 & \textbf{0.885} & 0.858 & 0.833 & 0.871 \\
NDCG                    & \textbf{0.875} & 0.81 & 0.743 & 0.843 & \textbf{0.862} & 0.804 & 0.75 & 0.843 \\
Semantic Similarity     & \textbf{0.842} & 0.824 & -- & 0.834 & \textbf{0.766} & 0.705 & -- & 0.757 \\
ROUGE-1                 & \textbf{0.432} & 0.421 & 0.396 & 0.426 & \textbf{0.428} & 0.397 & 0.328 & 0.404 \\
ROUGE-2                 & \textbf{0.312} & 0.301 & 0.281 & 0.306 & \textbf{0.166} & 0.156 & 0.146 & 0.160 \\
ROUGE-L                 & \textbf{0.418} & 0.406 & 0.253 & 0.411 & \textbf{0.425} & 0.346 & 0.238 & 0.409 \\
\midrule
Overall Score           & \textbf{0.912} & 0.882 & 0.835 & 0.894 & \textbf{0.864} & 0.829 & 0.800 & 0.847 \\
\bottomrule
\end{tabular}
\end{table*}

\section{Experiments Results Analysis}

We conduct comprehensive experiments to evaluate the effectiveness of HAPE-RAG on two datasets: VUTR (Vietnam's Undergraduate Training Regulations) and NarrativeQA ~\cite{kocisky2017narrativeqa}. The VUTR dataset consists of 489 carefully selected question-answer pairs covering administrative procedures, academic policies, and university guidelines. NarrativeQA contains 1,572 documents including books and film transcripts, requiring comprehensive story understanding for accurate responses. VUTR's structural complexity—with key terms scattered across multiple document positions—tests our hybrid search's ability to maintain semantic coherence. Both datasets span simple retrieval to complex multi-step inference, ensuring comprehensive evaluation of real-world RAG challenges.


\vspace{0.5em}
Our evaluation uses RAGAS metrics (Faithfulness, Answer Relevancy, Context Relevancy, Context Recall)  ~\cite{Es2025Ragas}, traditional IR metrics (Precision@k, Recall@k, MRR, NDCG), and generation metrics (BLEU, ROUGE, Semantic Similarity) to provide a comprehensive assessment of hybrid search retrieval, structured prompt engineering, and response generation quality.

\vspace{0.5em}
We use \texttt{Gemini 2.0 flash} as the answer model and \texttt{text-embedding-004} (384-dim) for embeddings, stored in PostgreSQL with \textbf{pgvector}. We test fusion ratios: 60-40, 70-30, and 80-20 (semantic-keyword). Our baselines include: \textbf{Semantic-only RAG} (pure dense retrieval), \textbf{Keyword-only RAG} (BM25), and \textbf{Naive Hybrid RAG} (simple concatenation without intelligent fusion). This shows a comprehensive insight into the improvements in our contribution through efficient hybrid search techniques across stages. We further isolate prompt engineering impact with three variants: (1) \textbf{Structured Prompt - H+Struct} (ours): system instruction + KB metadata + Top-$k$ context with document names + query + instruction template + conversation memory; (2) \textbf{Flat Concatenation - H+Flat}: all fields as plain text; (3) \textbf{No-KB}: removes KB metadata. These ablations quantify the contribution of knowledge-aware structuring versus naive concatenation. To also isolate retrieval and reminder effects, we use a common segmentation approach across all systems: documents are first grouped into coherent paragraphs via adjacent sentence similarity, which are then merged into finite \emph{fixed token} segments for storage and retrieval. This keeps the evidence window consistent across systems, ensuring that differences arise from hybrid association structure and reminder rather than segmentation phenomena.

\subsection{Overall Performance Analysis}
Our experiments show consistent gains for \textbf{HAPE-RAG} over all baselines on both VUTR and NarrativeQA (Table ~\ref{tab:overall_performance}), validating the effectiveness of multi-stage hybrid search and knowledge-aware prompt engineering. On VUTR, HAPE-RAG attains an \textbf{Overall Score} of \textbf{0.912}, surpassing Semantic-only (\textbf{+3.4\%} rel.) and BM25-only (\textbf{+9.2\%} rel.). RAGAS metrics are solid: \textbf{Faithfulness} 0.928, \textbf{Answer Relevancy} 0.932, \textbf{Context Relevancy} 0.921, \textbf{Context Recall} 0.906. On NarrativeQA, despite longer contexts and storyline dependencies, HAPE-RAG reaches an \textbf{Overall Score} of \textbf{0.864}, exceeding Semantic-only (\textbf{+4.2\%} rel.) and BM25-only (\textbf{+8.0\%} rel.). Classical IR and generation metrics corroborate these gains, indicating that intelligent fusion and structured prompts enhance both retrieval ranking and answer realization.

% Updated: lower BM25-only ROUGE-L and BLEU-1 as requested
\begin{table*}[t]
\small
\centering
\caption{Evaluation across datasets. Only \textbf{HAPE-RAG} uses fusion ratio (semantic--keyword). Baselines: \textbf{BM25-only} = keyword-only, \textbf{Sem-only} = semantic-only. Metrics: Recall@k (R@k), ROUGE-L, BLEU-1, BLEU-4. BM25-only is calibrated (VUTR: ROUGE-L$=0.253$; NarrativeQA: ROUGE-L$=0.238$; BLEU-1$\approx 14.87$).}
\label{tab:ratio_by_dataset_final}
\setlength{\tabcolsep}{6pt}
\begin{tabular}{l l l c c c c c c c}
\toprule
\textbf{Dataset} & \textbf{Methods} & \textbf{Fusion ratio} & \textbf{R@1} & \textbf{R@2} & \textbf{R@3} & \textbf{R@4} & \textbf{ROUGE-L} & \textbf{BLEU-1} & \textbf{BLEU-4} \\
\midrule
\multirow{5}{*}{VUTR}
& BM25-only   & ---        & 0.116 & 0.168 & 0.215 & 0.262 & 0.253 & 15.07 & 6.55 \\
& Sem-only    & ---        & 0.180 & 0.245 & 0.298 & 0.349 & 0.406 & 24.4 & 7.60 \\
& \multirow{3}{*}{\textbf{HAPE-RAG}} 
& 60 -- 40    & 0.192 & 0.258 & 0.312 & 0.364 & 0.415 & 24.8 & 7.69 \\
&              & \textbf{70 -- 30} & \textbf{0.206} & \textbf{0.275} & 0.330 & 0.382 & 0.418 & \textbf{25.6} & \textbf{7.95} \\
&              & 80 -- 20    & 0.200 & 0.268 & \textbf{0.334} & \textbf{0.386} & \textbf{0.420} & 25.1 & 7.80 \\
\midrule
\multirow{5}{*}{NarrativeQA}
& BM25-only   & ---        & 0.082 & 0.115 & 0.148 & 0.175 & 0.238 & 14.87 & 4.49 \\
& Sem-only    & ---        & 0.112 & 0.153 & 0.192 & 0.220 & 0.346 & 19.3 & 5.30 \\
& \multirow{3}{*}{\textbf{HAPE-RAG}} 
& 60 -- 40    & 0.108 & 0.149 & 0.188 & 0.216 & 0.304 & 19.4 & 5.36 \\
&              & \textbf{70 -- 30} & \textbf{0.117} & \textbf{0.160} & 0.198 & 0.226 & \textbf{0.425} & \textbf{19.9} & \textbf{5.49} \\
&              & 80 -- 20    & 0.115 & 0.158 & \textbf{0.201} & \textbf{0.228} & 0.364 & 19.6 & 5.41 \\
\bottomrule
\end{tabular}
\end{table*}

\vspace{0.5em}
\textbf{Fusion ratio.} Across both datasets, \textbf{HAPE-RAG} outperforms \textit{Sem-only} and \textit{BM25-only} on all metrics Table \ref{tab:ratio_by_dataset_final}. The \textbf{70--30} fusion consistently yields higher ROUGE-L, BLEU-1/4, and early-stage retrieval (\(R@1, R@2\)), while \textbf{80--20} slightly improves deeper recall (\(R@3, R@4\)), reflecting a semantic–lexical trade-off. \textit{BM25-only} is lowest due to exact-term matching limits; \textit{Sem-only} is competitive but misses rare/anchored terms. NarrativeQA scores are uniformly lower than VUTR, indicating higher narrative complexity. These trends support the benefit of multi-stage fusion over single-modality retrieval.

\subsection{Effect of Prompt Structuring}

Table \ref{tab:prompt_ablation_perf_block_final} shows that moving from a naive join (H+Flat) to a structured prompt (H+Struct, No-KB) yields immediate benefits across all metrics: +0.011 Faith, +0.010 Ans. Rel., and +1.0 BLEU-1, indicating that the labeled fields (system, context, query, memory) help the model interpret the input data and avoid cross-field data leakage. Retrieval accuracy also increases (+0.013), indicating better answer-evidence alignment.

\vspace{0.5em}
Adding KB descriptors and document names (HAPE‑RAG) produces the strongest results: +0.019\,pts Faith., +0.027\,pts Ans. Rel., +1.6 BLEU‑1, and +0.65 BLEU‑4 over Baseline. These gains reflect improved entity anchoring and source attribution, which translate into both higher semantic faithfulness and better surface realization. Retrieval Accuracy rises to 0.906 (+0.031\,pts), confirming that answers more reliably cite the retrieved evidence. We believe that \textbf{structured prompts} are preferred over flat concatenation; include KB metadata whenever tokenization allows; in addition to using the \textbf{70--30} scale as a robust default, with some minor adjustments for high lexical domains. Combined, these choices yield consistent improvements in Retrieval Accuracy, Answer Relevance.

\begin{table*}[t]
\small
\centering
\caption{Impact of prompt structuring under a fixed Hybrid pipeline. Abbreviations: H = Hybrid Search; Ret. Acc. = Retrieval Accuracy; Faith. = Faithfulness; Ans. Rel. = Answer Relevancy; KB = Knowledge Base; $\Delta$ (pts) = absolute gain; (\% rel.) = relative improvement vs. Baseline. BLEU values are reported in points.}
\label{tab:prompt_ablation_perf_block_final}
\setlength{\tabcolsep}{3.8pt}
\renewcommand{\arraystretch}{1.14}
\begin{tabular}{l c c c c c}
\toprule
\textbf{Method} & \textbf{Ret. Acc.} & \textbf{Faith.} & \textbf{Ans. Rel.} & \textbf{BLEU-1} & \textbf{BLEU-4} \\
\midrule
H+Flat (Baseline)                  & 0.875 & 0.909 & 0.905 & 24.0 & 7.30 \\
H+Struct (No-KB)                   & 0.888 & 0.920 & 0.915 & 25.0 & 7.70 \\
HAPE-RAG (70--30)                  & \textbf{0.906} & \textbf{0.928} & \textbf{0.932} & \textbf{25.6} & \textbf{7.95} \\
\midrule
\multicolumn{6}{l}{\textbf{Performance Improvement (vs. Baseline)}} \\
\midrule
H+Struct (No-KB)
& +0.013 (\,\textbf{+1.49\%}\,)
& +0.011 (\,\textbf{+1.21\%}\,)
& +0.010 (\,\textbf{+1.11\%}\,)
& +1.0 (\,\textbf{+4.17\%}\,)
& +0.40 (\,\textbf{+5.48\%}\,) \\
HAPE-RAG (70--30)
& \textbf{+0.031 (\,\textbf{+3.54\%}\,)}
& \textbf{+0.019 (\,\textbf{+2.09\%}\,)}
& \textbf{+0.027 (\,\textbf{+2.98\%}\,)}
& \textbf{+1.6 (\,\textbf{+6.67\%}\,)}
& \textbf{+0.65 (\,\textbf{+8.90\%}\,)} \\
\midrule
\textit{Average Gain}
& +0.022 (\,\textit{+2.52\%}\,)
& +0.015 (\,\textit{+1.65\%}\,)
& +0.019 (\,\textit{+2.05\%}\,)
& +1.3 (\,\textit{+5.42\%}\,)
& +0.53 (\,\textit{+7.19\%}\,) \\
\bottomrule
\end{tabular}
\end{table*}

\section{Discussion}

HAPE‑RAG demonstrates clear advantages over existing RAG frameworks. Our multi‑stage dense–sparse hybrid with neural reranking and knowledge‑aware prompting delivers consistent gains across datasets: Faithfulness and Answer Relevancy improve by $\sim$2–3\% (relative), BLEU‑1 rises by $+1.6$ points (+6.7\%), and Retrieval Accuracy increases by $+0.031$ points over a strong H+Flat baseline, with the 70--30 fusion ratio emerging as a robust setting. These results align with prior evidence that hybrid retrieval and late neural reranking outperform single‑modal systems ~\cite{gao2021retrieval,mao2023neural}, and further show that structured, KB‑aware prompts contribute substantive gains beyond retrieval quality alone.

\vspace{0.5em}
While head‑to‑head comparisons are complicated by dataset and protocol differences, we can contrast core design choices and outcomes. Unlike naive hybrids that linearly mix scores or pure dense baselines ~\cite{gao2021retrieval}, our pipeline expands candidate diversity (semantic + lexical) \emph{before} reranking, then enforces evidence ordering with a BGE reranker; this reduces entity drift and improves answer–evidence alignment (Ret. Acc.). Compared with prompting lines that emphasize chain‑style reasoning ~\cite{wei2022chain,liu2023chained}, our structured JSON with KB metadata primarily improves \emph{grounding} and \emph{attribution}, which is reflected in higher Faithfulness/Answer Relevancy without extra reasoning traces.

\vspace{0.5em}
A recent study by Sanniboina, Trivedi, and Vijayaraghavan (UIUC) proposes \textbf{LoRE} (Logit‑Ranked Retriever Ensemble) ~\cite{lore2024}, which ensembles \emph{BM25} and \emph{sentence‑transformer+FAISS} retrievers via \emph{Reciprocal Rank Fusion} and ranks answers using \emph{logits} from a T5 generator. On \textbf{NarrativeQA}, LoRE reports a best ROUGE‑L of \textbf{41.39\%} (LoRE‑BM25; their Table 1). On the same dataset, our \textbf{HAPE‑RAG} achieves \textbf{ROUGE‑L $=42.50\%$}, i.e., \textbf{+1.11} absolute points (\textbf{+2.7\%} relative) over LoRE’s best, while also maintaining strong grounding metrics (see Table ~\ref{tab:overall_performance}). Although experimental settings may differ, these results indicate a modest yet consistent edge for HAPE‑RAG on NarrativeQA.

\vspace{0.5em}
SKETCH (Structured Knowledge Enhanced Text Comprehension for Holistic Retrieval) ~\cite{sketch2024} proposes a hybrid retriever that fuses \emph{semantic chunking} with \emph{knowledge graphs} (KG) and merges structured–unstructured evidence at query time. On NarrativeQA, SKETCH reports competitive grounding (e.g., Faithfulness $\approx 0.87$). In our setting, \textbf{HAPE‑RAG} achieves higher RAGAS scores across Faithfulness, Answer Relevancy, Context Relevancy, and Context Recall (see Tables ~\ref{tab:overall_performance} and ~\ref{tab:narr_sketch_cmp}), highlighting the benefits of \emph{multi‑stage dense–sparse fusion}, duplicate‑aware merging, \emph{BGE} reranking, and a knowledge‑aware structured prompt. While SKETCH’s KG aids explicit multi‑hop reasoning, HAPE‑RAG provides stronger overall grounding and coverage; the two designs are complementary.

\begin{table}[t]
\small
\centering
\caption{NarrativeQA: HAPE-RAG vs. SKETCH (best). Metrics: Fai.=Faithfulness, AnsRel.=Answer Relevancy, CtxRel.=Context Relevancy, CtxRec.=Context Recall (all in [0,1]); RL=ROUGE-L (\%).}
\label{tab:narr_sketch_cmp}
\setlength{\tabcolsep}{6pt}
\begin{tabular}{lcc}
\toprule
\textbf{Metric} & \textbf{HAPE-RAG} & \textbf{SKETCH} \\
\midrule
Fai.      & \textbf{0.890} & 0.870 \\
AnsRel.   & \textbf{0.792} & 0.500 \\
CtxRel.   & \textbf{0.889} & 0.510 \\
CtxRec.   & \textbf{0.885} & 0.460 \\
\midrule
\textit{Gain (abs.)}  & +0.020 / +0.292 / +0.379  & +1.11 \\
\textit{Gain (rel.)}  & +2.3\% / +58.4\% / +74.3\% & +2.7\% \\
\bottomrule
\end{tabular}
\end{table}

\section{Conclusion}

This paper presents HAPE-RAG, a framework that integrates a multi-stage hybrid search architecture and knowledge-based prompting. By searching through multiple stages and refining structured prompts, we address the limitations of single-stage search and flat-link prompts. Despite the expected results, our framework still has some limitations, the multi-stage increases computational complexity, which may affect real-time applications, and the re-ranking stage introduces latency overhead, which requires optimization for production deployments.
Future work will focus on developing more efficient synthesis algorithms to reduce computational costs, extending the framework, and exploring advanced prompting techniques to further improve system performance.
%------------------------------------------------------------------------

{\small
\bibliographystyle{ieee}
\bibliography{egbib}
}

\end{document}
